% GENERATED FILE. Edit canonical lessons, then run npm run generate:books.
% canonical-source-hash: fnv1a64:f0c0f146be555f05
% canonical-lessons: GU-C516-cuno, GU-C516-dholvu-whitewash, GU-C516-rang, GU-C516-rangvu, GU-C516-lipvu, GU-R516-review-509-516

\chapter{Lime and Paint}
\label{ch:gu-lime-and-paint}
\hlchaptermodality{516}

\begin{chapteropening}
\textbf{By the end of this chapter:} \emph{I can talk about whitewashing and painting a house.}

Complete the last lesson of chapter 516: I can talk about whitewashing and painting a house.
\end{chapteropening}

\section[\texorpdfstring{\gu{ચૂનો} (cūno)}{ચૂનો (cūno)}]{\gu{ચૂનો} (cūno) — lime wash}

\label{lesson:GU-C516-cuno}



\begin{quote}
Before the new one: say the Gujarati for a pipe, then the Gujarati for a plumber.
\end{quote}

\subsection*{You'll want to know: \gu{ચૂનો}}
\textbf{\gu{ચૂનો}} — \emph{cūno} — \textquotedblleft{}lime wash\textquotedblright{}.

\begin{quote}
\gu{દિવાળી} \gu{પહેલાં} \gu{અમે} \gu{દીવાલો} \gu{પર} \gu{ચૂનો} \gu{લગાડીએ} \gu{છીએ}.
\end{quote}

\begin{quote}
\emph{divāḷī pahelā̃ ame dīvālo par cūno lagāḍīe chīe.}
\end{quote}

\begin{quote}
Before Diwali we put lime wash on the walls.
\end{quote}

A masculine noun.

\begin{cousinweb}
\textbf{\gu{ચૂનો}} (\emph{cūno}) is a tadbhava, inherited from Sanskrit \emph{cūrṇa}, \textquotedblleft{}powder\textquotedblright{}, and CDIAL lists the Gujarati word as \textquotedblleft{}lime\textquotedblright{} (CDIAL 4889).
\end{cousinweb}

\begin{grammarlens}[title={What you've built}]
One more word for describing what you lived and felt.
\end{grammarlens}

\subsection*{Your turn}
\begin{itemize}
\raggedright
  \item \emph{Say it:} \emph{cūno}
  \item \emph{Say it:} \emph{cūno}, once more
  \item \emph{Say it:} \emph{divāḷī pahelā̃ ame dīvālo par cūno lagāḍīe chīe.} — Before Diwali we put lime wash on the walls.
  \item \emph{Recall:} say the Gujarati for a pipe, then the Gujarati for a plumber, then say \emph{cūno} again
\end{itemize}

\begin{tcolorbox}[breakable,colback=teal!4,colframe=teal!35!black,title={Before you move on}]
What does \gu{ચૂનો} mean? (\textquotedblleft{}Lime wash\textquotedblright{}.) Say it once more.
\end{tcolorbox}
\section[\texorpdfstring{\gu{ધોળવું} (dhoḷvũ)}{ધોળવું (dhoḷvũ)}]{\gu{ધોળવું} (dhoḷvũ) — to whitewash}

\label{lesson:GU-C516-dholvu-whitewash}



\begin{quote}
Before the new one: say the Gujarati for a plumber, then the Gujarati for lime wash.
\end{quote}

\subsection*{You'll want to know: \gu{ધોળવું}}
\textbf{\gu{ધોળવું}} — \emph{dhoḷvũ} — \textquotedblleft{}to whitewash\textquotedblright{}.

\begin{quote}
\gu{ગઈકાલે} \gu{અમે} \gu{આખું} \gu{ઘર} \gu{ધોળ્યું}.
\end{quote}

\begin{quote}
\emph{gaīkāle ame ākhũ ghar dhoḷyũ.}
\end{quote}

\begin{quote}
Yesterday we whitewashed the whole house.
\end{quote}

Past, with \gu{તેણે} (he or she …): \gu{તેણે} \gu{ઘર} \gu{ધોળ્યું}, \emph{teṇe ghar dhoḷyũ}.

\begin{cousinweb}
\textbf{\gu{ધોળવું}} (\emph{dhoḷvũ}) is inherited from Sanskrit \emph{dhavalá}, \textquotedblleft{}white\textquotedblright{}, and CDIAL lists the Gujarati verb, \textquotedblleft{}to whitewash\textquotedblright{} (CDIAL 6767).
\end{cousinweb}

\begin{grammarlens}[title={What you've built}]
One more word for describing what you lived and felt.
\end{grammarlens}

\subsection*{Your turn}
\begin{itemize}
\raggedright
  \item \emph{Say it:} \emph{dhoḷvũ}
  \item \emph{Say it:} \emph{dhoḷvũ}, once more
  \item \emph{Say it:} \emph{gaīkāle ame ākhũ ghar dhoḷyũ.} — Yesterday we whitewashed the whole house.
  \item \emph{Recall:} say the Gujarati for a plumber, then the Gujarati for lime wash, then say \emph{dhoḷvũ} again
\end{itemize}

\begin{tcolorbox}[breakable,colback=teal!4,colframe=teal!35!black,title={Before you move on}]
What does \gu{ધોળવું} mean? (\textquotedblleft{}To whitewash\textquotedblright{}.) Say it once more.
\end{tcolorbox}
\section[\texorpdfstring{\gu{રંગ} (rang)}{રંગ (rang)}]{\gu{રંગ} (rang) — a colour}

\label{lesson:GU-C516-rang}



\begin{quote}
Before the new one: say the Gujarati for lime wash, then the Gujarati for to whitewash.
\end{quote}

\subsection*{You'll want to know: \gu{રંગ}}
\textbf{\gu{રંગ}} — \emph{rang} — \textquotedblleft{}a colour\textquotedblright{}.

\begin{quote}
\gu{તમને} \gu{કયો} \gu{રંગ} \gu{ગમે} \gu{છે}?
\end{quote}

\begin{quote}
\emph{tamne kayo rang game che?}
\end{quote}

\begin{quote}
Which colour do you like?
\end{quote}

A masculine noun.

\begin{cousinweb}
\textbf{\gu{રંગ}} (\emph{rang}) is a tatsama, taken from Sanskrit \emph{raṅga}, \textquotedblleft{}colour, paint, dye\textquotedblright{} (Monier-Williams). It also means paint: \textbf{\gu{લીલો} \gu{રંગ}}, green paint.
\end{cousinweb}

\begin{grammarlens}[title={What you've built}]
One more word for describing what you lived and felt.
\end{grammarlens}

\subsection*{Your turn}
\begin{itemize}
\raggedright
  \item \emph{Say it:} \emph{rang}
  \item \emph{Say it:} \emph{rang}, once more
  \item \emph{Say it:} \emph{tamne kayo rang game che?} — Which colour do you like?
  \item \emph{Recall:} say the Gujarati for lime wash, then the Gujarati for to whitewash, then say \emph{rang} again
\end{itemize}

\begin{tcolorbox}[breakable,colback=teal!4,colframe=teal!35!black,title={Before you move on}]
What does \gu{રંગ} mean? (\textquotedblleft{}A colour\textquotedblright{}.) Say it once more.
\end{tcolorbox}
\section[\texorpdfstring{\gu{રંગવું} (rangvũ)}{રંગવું (rangvũ)}]{\gu{રંગવું} (rangvũ) — to colour (with paint or dye)}

\label{lesson:GU-C516-rangvu}



\begin{quote}
Before the new one: say the Gujarati for to whitewash, then the Gujarati for a colour.
\end{quote}

\subsection*{You'll want to know: \gu{રંગવું}}
\textbf{\gu{રંગવું}} — \emph{rangvũ} — \textquotedblleft{}to colour (with paint or dye)\textquotedblright{}.

\begin{quote}
\gu{અરુણે} \gu{બારણું} \gu{લીલા} \gu{રંગથી} \gu{રંગ્યું}.
\end{quote}

\begin{quote}
\emph{Aruṇe bārṇũ līlā rangthī rangyũ.}
\end{quote}

\begin{quote}
Arun painted the door green.
\end{quote}

Past, with \gu{તેણે} (he or she …): \gu{તેણે} \gu{બારણું} \gu{રંગ્યું}, \emph{teṇe bārṇũ rangyũ}.

\begin{cousinweb}
\textbf{\gu{રંગવું}} (\emph{rangvũ}) is built on \textbf{\gu{રંગ}}, \textquotedblleft{}colour\textquotedblright{}, from Sanskrit \emph{raṅga}, \textquotedblleft{}colour, dye\textquotedblright{} (Monier-Williams); it is also used for dyeing cloth.
\end{cousinweb}

\begin{grammarlens}[title={What you've built}]
One more word for describing what you lived and felt.
\end{grammarlens}

\subsection*{Your turn}
\begin{itemize}
\raggedright
  \item \emph{Say it:} \emph{rangvũ}
  \item \emph{Say it:} \emph{rangvũ}, once more
  \item \emph{Say it:} \emph{Aruṇe bārṇũ līlā rangthī rangyũ.} — Arun painted the door green.
  \item \emph{Recall:} say the Gujarati for to whitewash, then the Gujarati for a colour, then say \emph{rangvũ} again
\end{itemize}

\begin{tcolorbox}[breakable,colback=teal!4,colframe=teal!35!black,title={Before you move on}]
What does \gu{રંગવું} mean? (\textquotedblleft{}To colour (with paint or dye)\textquotedblright{}.) Say it once more.
\end{tcolorbox}
\section[\texorpdfstring{\gu{લીંપવું} (lī̃pvũ)}{લીંપવું (lī̃pvũ)}]{\gu{લીંપવું} (lī̃pvũ) — to plaster (with mud)}

\label{lesson:GU-C516-lipvu}



\begin{quote}
Before the new one: say the Gujarati for a colour, then the Gujarati for to colour.
\end{quote}

\subsection*{You'll want to know: \gu{લીંપવું}}
\textbf{\gu{લીંપવું}} — \emph{lī̃pvũ} — \textquotedblleft{}to plaster (with mud)\textquotedblright{}.

\begin{quote}
\gu{ગામમાં} \gu{દાદી} \gu{ગારાથી} \gu{આંગણું} \gu{લીંપતી} \gu{હતી}.
\end{quote}

\begin{quote}
\emph{gāmmā̃ dādī gārāthī ā̃gaṇũ lī̃ptī hatī.}
\end{quote}

\begin{quote}
In the village, my grandmother used to plaster the courtyard with mud.
\end{quote}

Past, with \gu{તેણે} (he or she …): \gu{તેણે} \gu{આંગણું} \gu{લીંપ્યું}, \emph{teṇe ā̃gaṇũ lī̃pyũ}.

\begin{cousinweb}
\textbf{\gu{લીંપવું}} (\emph{lī̃pvũ}) is inherited from Sanskrit \emph{limpáti}, \textquotedblleft{}smears\textquotedblright{}, and CDIAL lists the Gujarati verb (CDIAL 11066).
\end{cousinweb}

\begin{grammarlens}[title={What you've built}]
One more word for describing what you lived and felt.
\end{grammarlens}

\subsection*{Your turn}
\begin{itemize}
\raggedright
  \item \emph{Say it:} \emph{lī̃pvũ}
  \item \emph{Say it:} \emph{lī̃pvũ}, once more
  \item \emph{Say it:} \emph{gāmmā̃ dādī gārāthī ā̃gaṇũ lī̃ptī hatī.} — In the village, my grandmother used to plaster the courtyard with mud.
  \item \emph{Recall:} say the Gujarati for a colour, then the Gujarati for to colour, then say \emph{lī̃pvũ} again
\end{itemize}

\begin{tcolorbox}[breakable,colback=teal!4,colframe=teal!35!black,title={Before you move on}]
What does \gu{લીંપવું} mean? (\textquotedblleft{}To plaster (with mud)\textquotedblright{}.) Say it once more.
\end{tcolorbox}
\section[(dialogue)]{Looking back: from Bellows, Hammer and Rust to Lime and Paint}

\label{lesson:GU-R516-review-509-516}



\begin{quote}
Say the Gujarati for bellows and to blow out.
\end{quote}

\subsection*{Your turn}
Say these aloud:

\begin{itemize}
\raggedright
  \item bellows, to blow out, to beat (metal) flat, rust, to rust
  \item a goldsmith, to melt down, copper, to cast, to have (something) made
  \item a labourer, a foreman, a mason, a mechanic, a clerk
  \item a job post, a job interview, a wage, a career, a promotion
  \item an industry, a company, a machine, brass, scrap
  \item loose, a screw, to tighten, a wall peg, to hang (something)
  \item a flake (of plaster), to peel off, a hole, a pipe, a plumber
  \item lime wash, to whitewash, a colour, to colour, to plaster
\end{itemize}

\begin{tcolorbox}[breakable,colback=teal!4,colframe=teal!35!black,title={Before you move on}]
Bellows? (\textbf{\gu{ધમણ}}, \emph{dhamaṇ}.) To blow out? (\textbf{\gu{ફૂંકવું}}, \emph{phū̃kvũ}.) To beat (metal) flat? (\textbf{\gu{ટીપવું}}, \emph{ṭīpvũ}.) Rust? (\textbf{\gu{કાટ}}, \emph{kāṭ}.) To rust? (\textbf{\gu{કટાવું}}, \emph{kaṭāvũ}.) A goldsmith? (\textbf{\gu{સોની}}, \emph{sonī}.) To melt down? (\textbf{\gu{ઓગાળવું}}, \emph{ogāḷvũ}.) Copper? (\textbf{\gu{તાંબું}}, \emph{tā̃bũ}.) To cast? (\textbf{\gu{ઢાળવું}}, \emph{ḍhāḷvũ}.) To have (something) made? (\textbf{\gu{ઘડાવવું}}, \emph{ghaḍāvvũ}.) A labourer? (\textbf{\gu{મજૂર}}, \emph{majūr}.) A foreman? (\textbf{\gu{મુકાદમ}}, \emph{mukādam}.) A mason? (\textbf{\gu{કડિયો}}, \emph{kaḍiyo}.) A mechanic? (\textbf{\gu{મિસ્ત્રી}}, \emph{mistrī}.) A clerk? (\textbf{\gu{કારકુન}}, \emph{kārkun}.) A job post? (\textbf{\gu{પદ}}, \emph{pad}.) A job interview? (\textbf{\gu{ઇન્ટરવ્યૂ}}, \emph{inṭaravyū}.) A wage? (\textbf{\gu{વેતન}}, \emph{vetan}.) A career? (\textbf{\gu{કારકિર્દી}}, \emph{kārkirdī}.) A promotion? (\textbf{\gu{બઢતી}}, \emph{baḍhtī}.) An industry? (\textbf{\gu{ઉદ્યોગ}}, \emph{udyog}.) A company? (\textbf{\gu{કંપની}}, \emph{kampnī}.) A machine? (\textbf{\gu{મશીન}}, \emph{maśīn}.) Brass? (\textbf{\gu{પિત્તળ}}, \emph{pittaḷ}.) Scrap? (\textbf{\gu{ભંગાર}}, \emph{bhangār}.) Loose? (\textbf{\gu{ઢીલું}}, \emph{ḍhīlũ}.) A screw? (\textbf{\gu{સ્ક્રૂ}}, \emph{skrū}.) To tighten? (\textbf{\gu{કસવું}}, \emph{kasvũ}.) A wall peg? (\textbf{\gu{ખીંટી}}, \emph{khī̃ṭī}.) To hang (something)? (\textbf{\gu{લટકાવવું}}, \emph{laṭkāvvũ}.) A flake (of plaster)? (\textbf{\gu{પોપડો}}, \emph{popḍo}.) To peel off? (\textbf{\gu{ઉખડવું}}, \emph{ukhaḍvũ}.) A hole? (\textbf{\gu{કાણું}}, \emph{kāṇũ}.) A pipe? (\textbf{\gu{પાઈપ}}, \emph{pāīp}.) A plumber? (\textbf{\gu{પ્લમ્બર}}, \emph{plambar}.) Lime wash? (\textbf{\gu{ચૂનો}}, \emph{cūno}.) To whitewash? (\textbf{\gu{ધોળવું}}, \emph{dhoḷvũ}.) A colour? (\textbf{\gu{રંગ}}, \emph{rang}.) To colour? (\textbf{\gu{રંગવું}}, \emph{rangvũ}.) To plaster? (\textbf{\gu{લીંપવું}}, \emph{lī̃pvũ}.)
\end{tcolorbox}
